\section{Introduction}\label{sec:intro}

A translation shown beside the text it translates asks the reader to look twice: once at the original, once at the
answer. A translation painted \emph{over} the original asks nothing, but it must be right in a stricter sense. The patch has to
cover exactly the old glyphs and nothing else, the new text has to fit the room the old text had, it has to look like the
page it sits on, and it has to stay put when the page moves. Mwoyeo\footnote{All code in this paper is copied from the
repository at revision \texttt{499e312} (2026-10-19), the main branch when the paper was written. Paths are relative to the
repository root; \texttt{sources.tsv} lists, for every number and every claim about the code, the file and line that carries
it, and \texttt{check.py} fails when one of them has moved.} is a macOS application that does this for the area under a
movable \emph{lens}.

The last requirement is the one that separates a demonstration from a tool. A lens over a chat or a feed is used while the
person scrolls it. The capture, the reading of the text, the translation and the repainting take time, so the translations
that arrive describe a frame that is already gone. Left alone they trail the text they translate by the speed of the scroll
times the age of the frame; the eye reads that as a shimmer, and in the worst case a translation sits on the wrong line.
That latency matters for direct manipulation is well established \citep{latency}.

\paragraph{Contributions.}
\begin{itemize}
\item A scroll predictor that places translations where the content will be by the time they are shown, from wheel and
  trackpad events alone, and that learns how far one wheel unit moves the page in the application under the pointer
  (Section~\ref{sec:scroll}). It is written as numbered equations, an algorithm and a 38-line listing, with a drift
  measure that the self-test prints.
\item A rule set for what must \emph{not} be translated: a formula or an identifier line, prose that is mostly inline mathematics,
  and regions of code, terminal output, logs, configuration and diffs found from tint, indentation and line-number gutters
  (Section~\ref{sec:keep}). The rules read text and geometry only; no model, no word list per language.
\item An erasure and re-setting stage built on a robust estimate of the background, with MI-GAN \citep{migan} where the
  background is not flat (Section~\ref{sec:erase}).
\item A description of how each part is measured in the repository's benches, with numbers taken only from result files
  that are committed there (Section~\ref{sec:measure}).
\end{itemize}

\paragraph{Scope.}
The paper describes what the code does at the stated revision, including what it does badly (Section~\ref{sec:limits}). Text
recognition runs in two places in this revision: document reads run in an XPC service, line reads still run in the
application's own process; moving the line reads into a second service is in progress and is not described as done.
We do not evaluate translation quality in general: the translation engines are interchangeable components, and the benches
here measure what the \emph{screen} adds to the problem.
